% Image credits for the photographs and AI illustrations of Book 2 (High School).
% Machine-readable license records live in images/book2/CREDITS.md; keep the
% two files in sync. The Book 2 writer owns this file: add one \item per
% photograph (inside the itemize, by chapter) as photographs land.
\chapter*{图片版权说明}

本书中的插图、地图和图表均为原创；依据已发表数据绘制的地图和图表，在图注或本列表中注明来源。照片依各自的自由许可协议使用，谨向其作者致谢；每张照片的完整来源链接和许可协议网址，均列于本书代码库中的 \texttt{images/book2/CREDITS.md}。

\bigskip
除照片和原创插图外，部分图为 AI 生成的插图，由 OpenAI 图像生成功能根据本书编辑撰写的提示词生成，并在收录前经过科学准确性审阅。这些图仅展示通用场景，从不呈现真实地点、真实人物或真实标本。每张生成图像的完整提示词，均记录于代码库中的 \texttt{images/book2/ai/PROMPTS.md}。

\bigskip
\noindent 照片与数据，按章列出：
\begin{itemize}\small
  \item 第~1 章（测量地球）：阿波罗~17 号拍摄的地球；照片由 NASA 提供，公有领域（维基共享资源）。
  \item 第~1 章：一幅可重复的圆形插图，取自 J.-B. Delambre，《Base du système métrique décimal》，第~2 卷（1807）；公有领域（维基共享资源，来自法国国家图书馆 Gallica）。
  \item 第~1 章：经度与纬度每度长度，依据 WGS~84 椭球体计算（EPSG 数据集）。
  \item 第~2 章（宜居行星）：水星号探测器（Mariner~10）拍摄的金星（PIA23791），NASA/JPL-Caltech，公有领域；好奇号火星车拍摄的夏普山（PIA16769），NASA/JPL-Caltech/MSSS，公有领域（维基共享资源）。
  \item 第~2 章：行星数据取自 NASA 概况表；水的相界取自 IAPWS 发布的 SR1-86 和 R14-08；宜居带界限取自 Kopparapu 等人（2013）；比邻星 b 取自 NASA 系外行星档案。
  \item 第~3 章（地球的能量收支）：北极海冰上的融池，摄于 2011 年 ICESCAPE 任务期间；照片由 NASA 戈达德太空飞行中心提供，公有领域（维基共享资源）。
  \item 第~3 章：能量收支数值取自 IPCC 第六次评估报告第一工作组第~7.2 图；反照率取自科罗拉多大学博尔德分校（MOSAiC）和美国国家冰雪数据中心；热量输送取自 Trenberth 与 Caron（2001）；净初级生产力取自 Field 等人（1998）；热化学数据取自 NIST 化学网络手册。
  \item 第~4 章（云为何形成）：晴天积云，照片由拉尔夫·F.·克雷奇拍摄，NOAA 提供；积雨云的砧状云，NOAA 国家强风暴实验室提供；两者均为公有领域（维基共享资源）。
  \item 第~4 章：饱和度依据 IAPWS 发布的数据计算；水滴尺寸取自 NASA 喷气推进实验室。
  \item 第~5 章（风暴与热带气旋）：国际空间站拍摄的飓风伊莎贝尔，照片由迈克·特伦查德拍摄，NASA 约翰逊太空中心提供；堪萨斯州的龙卷风，照片由肖恩·沃夫拍摄，NOAA 国家强风暴实验室提供；闪电，照片由 C.~克拉克拍摄，NOAA 国家强风暴实验室提供；均为公有领域（维基共享资源）。
  \item 第~5 章：气旋路径取自 IBTrACS 第~4r01 版（Gahtan 等人，2024；Knapp 等人，2010；NOAA 国家环境信息中心）；飓风等级取自 NOAA 国家飓风中心；飓风能量与形成条件取自 NOAA 大西洋海洋与气象实验室；闪电频率取自 Christian 等人（2003）；龙卷风数据取自 NOAA 风暴预报中心。
  \item 第~6 章（分层的海洋）：采样转盘，照片由巴勃罗·克莱门特-科隆拍摄，NOAA 照片库提供；从科考船放出的剖面浮标，NOAA 全球海洋监测与观测计划提供；2002--2010 年海洋叶绿素，杰西·艾伦与罗伯特·西蒙，NASA 地球观测站；均为公有领域（维基共享资源）。
  \item 第~6 章：海洋剖面取自《世界海洋图集 2023》（Reagan 等人，2024；NOAA 国家环境信息中心）；海水密度取自国际状态方程（Fofonoff 与 Millard，1983）；海洋体积取自 Charette 与 Smith（2010）；浮标数据取自 Argo 项目办公室；上升流取自 NOAA 国家海洋局和 Kessler（2006）。
  \item 第~7 章（土壤，一种形成缓慢的资源）：喀拉拉邦片麻岩上的红土剖面，照片由维尔纳·舍尔曼拍摄，CC~BY-SA~2.5；马达加斯加中部的风化层，照片由 Diorit 拍摄，CC~BY-SA~3.0；土壤剖面，美国自然资源保护局提供，公有领域（均为维基共享资源）。
  \item 第~7 章：土壤生成与侵蚀速率取自 Montgomery（2007）；矿物化学式取自矿物数据库（webmineral.com）；风化层厚度取自美国地质调查局《美国地下水地图集》。
  \item 第~8 章（流域、河流与洪水）：测量河流流量的水文学家和一个测量站，美国地质调查局提供，公有领域（维基共享资源）。
  \item 第~8 章：华盛顿附近波托马克河的年度洪峰与日流量，取自美国地质调查局国家水信息系统，站点编号 01646500；密西西比河流量与流域取自 Kammerer（1990，USGS）；亚马孙河流量取自 Callède 等人（2010）；洪水安全信息取自 NOAA 国家气象局。
  \item 第~9 章（淡水资源）：2000 年和 2018 年的咸海，NASA 地球观测站（MODIS/Terra）；堪萨斯州萨布莱特附近的灌溉农田，NASA（ASTER/Terra）提供，公有领域（维基共享资源）。
  \item 第~9 章：全球水储量依据 Shiklomanov（1993），取自美国地质调查局水科学学校；全球水循环依据 Trenberth 等人（2007），转引自 Trenberth 与 Asrar（2014）；努比亚地下水年龄取自 Sturchio 等人（2004）；高平原含水层变化取自 McGuire（2017，USGS）；咸海水量收支取自 Gaybullaev 等人（2012）和 Micklin（2007）。
  \item 第~10 章（从大陆漂移到海底扩张）：钻探船 \emph{JOIDES Resolution}，大洋钻探计划，得克萨斯农工大学，公有领域（维基共享资源）。
  \item 第~10 章：Cande 与 Kent（1995）的极性时间尺度；深海钻探计划（1970）第~3 航次的钻探结果；海底年龄取自 Seton 等人（2020），EarthByte，许可协议 CC~BY~4.0。
  \item 第~11 章（地震波揭示地球内部）：在波多黎各安装临时地震仪；照片由托马斯·普拉特拍摄，美国地质调查局提供，公有领域（维基共享资源）。
  \item 第~11 章：Dziewonski 与 Anderson（1981）的《初步参考地球模型》；P 波与 S 波的经验走时取自澳大利亚被动地震服务器（AuSPASS）。
  \item 第~12 章（岩石圈与软流圈）：一座冰山，NOAA 提供，公有领域（取自美国地质调查局水科学学校）。
  \item 第~12 章：热流取自 Davies 与 Davies（2010），许可协议 CC~BY~3.0；热流区取自 Roy、Blackwell 与 Decker（1972）；地幔固相线取自 Hirschmann（2000）；地壳厚度与密度取自 Christensen 与 Mooney（1995）；高程取自 ETOPO1（NOAA）。
  \item 第~13 章（板块及其运动）：喀斯喀特山脉的一处卫星定位站，照片由戴夫·舍罗德拍摄，美国地质调查局提供，公有领域（维基共享资源）。北太平洋海底的晕渲地形图，美国国家地球物理数据中心，NOAA 提供，公有领域（维基共享资源）。
  \item 第~13 章：板块旋转极 NNR-MORVEL56 取自 Argus、Gordon 与 DeMets（2011）；台站位置取自内华达大地测量实验室（Blewitt、Hammond 与 Kreemer，2018）；火山年龄取自 Clague 与 Dalrymple（1987），美国地质调查局；板块边界取自 Bird（2003）；陆地数据取自 Natural Earth。
  \item 第~14 章（洋脊）：加拉帕戈斯裂谷上的枕状熔岩，NOAA Okeanos 探索计划提供，公有领域；胡安·德富卡洋脊上的黑烟囱，华盛顿大学与 NOAA 提供，公有领域；辉长岩，照片由达拉·桑德罗尔（GeoDIL）拍摄，CC0；玄武岩中的橄榄岩，照片由詹姆斯·圣约翰拍摄，CC~BY~2.0（均为维基共享资源）。
  \item 第~14 章：海底深度取自 ETOPO1（NOAA），海底年龄取自 Seton 等人（2020），EarthByte，许可协议 CC~BY~4.0；深度定律取自 Parsons 与 Sclater（1977）；地幔固相线取自 Hirschmann（2000）；旋转极取自 Argus、Gordon 与 DeMets（2011）的 NNR-MORVEL56。
  \item 第~15 章（俯冲带）：蓝片岩与榴辉岩，照片由詹姆斯·圣约翰拍摄，CC~BY~2.0；2011 年日本气仙沼的海啸破坏，照片由戴伦·麦科德拍摄，美国海军提供，公有领域（均为维基共享资源）。
  \item 第~15 章：地震数据取自美国地质调查局地震目录（ComCat）；板块边界与汇聚速率取自 Bird（2003）；海底年龄取自 Seton 等人（2020），EarthByte，许可协议 CC~BY~4.0；板块速度取自 Argus、Gordon 与 DeMets（2011）的 NNR-MORVEL56；陆地数据取自 Natural Earth。
  \item 第~16 章（矿物）：石英、正长石、黑云母和橄榄石，照片分别由沙农·海因勒、达拉·桑德罗尔、德克斯特·珀金斯和内萨·尤尔拍摄（地质数字图像库），CC0；薄片中的玄武岩，照片由 Chmee2 拍摄，CC~BY-SA~3.0；花岗岩薄片中的石英，照片由迈克尔·C.·里格尔拍摄，CC~BY-SA~3.0（均为维基共享资源）。
  \item 第~16 章：大陆地壳成分取自 Rudnick 与 Gao（2014）；标准原子量取自 IUPAC/CIAAW；矿物化学式与密度取自 webmineral.com；内华达山脉花岗质岩石的密度取自美国地质调查局（开放文件报告 91-570-A）。
  \item 第~17 章（火成岩）：巨人堤的玄武岩柱，照片由亚当·毕晓普拍摄，CC~BY-SA~3.0；优胜美地国家公园的酋长岩，美国地质调查局，CC0；玄武岩、安山岩、流纹岩、闪长岩、花岗岩、黑曜石、斑状岩及两张薄片，照片分别由内萨·尤尔、特维拉·贝克、约瑟夫·H.·哈特曼、迈克·戴维斯、德克斯特·珀金斯和达拉·桑德罗尔拍摄（地质数字图像库），CC0（均为维基共享资源）。
  \item 第~17 章：岩墙与岩床的冷却依据 Delaney（1986），美国地质调查局开放文件报告 86-444；岩浆成分与结晶顺序取自 S.~A.~Nelson 的《自然地质学》讲义和 S.~Earle 的《Physical Geology》（CC~BY~4.0）；鲍温的生平取自卡内基科学院地球物理实验室的历史资料。
  \item 第~18 章（煤、石油与天然气）：富碳页岩中的化石植物、褐煤、烟煤和无烟煤，照片由詹姆斯·圣约翰拍摄，CC~BY~2.0；罗马尼亚的一个褐煤层，美国地质调查局提供，公有领域（均为维基共享资源）。
  \item 第~18 章：油窗、煤的形成与圈闭取自 S.~Earle 的《Physical Geology》（CC~BY~4.0）；煤的碳含量取自美国能源信息署；石油系统依据 Magoon 与 Schmoker（美国地质调查局，2000）；中非西部的裂谷与盐盆取自美国地质调查局概况表 2023-3028。
  \item 第~19 章（大陆地壳）：阿卡斯塔片麻岩，照片由佩德罗亚历山德拉德拍摄，CC~BY-SA~3.0；杰克山锆石，图像由 A.~卡沃西和 J.~瓦利拍摄，公有领域；加拿大地盾的片麻岩和花岗岩，照片由克里斯·赖特拍摄，CC~BY-SA~4.0（均为维基共享资源）。
  \item 第~19 章：大陆地壳成分取自 Rudnick 与 Gao（2014）；大陆地壳厚度取自 Christensen 与 Mooney（1995）、White、McKenzie 与 O'Nions（1992）以及 Karplus 等人（2019）；大陆范围取自 Cogley（1984）；海底年龄取自 Seton 等人（2020），EarthByte，许可协议 CC~BY~4.0。
  \item 第~20 章（造山运动）：从哥伦比亚号航天飞机上拍摄的珠穆朗玛峰，NASA 提供，公有领域；卡利甘达基峡谷中的褶皱岩层，照片由让-马里·于洛拍摄，CC~BY~2.0；格拉鲁斯逆冲断层，照片由克里斯蒂安·海涅拍摄，CC~BY-SA~3.0；塞迈勒蛇绿岩的枕状熔岩，照片由利奥内尔·迪乔伊泽勒拍摄，CC~BY-SA~3.0（均为维基共享资源）。
  \item 第~20 章：印度和拉萨的过去位置取自 Müller 等人（2019、2022）的板块模型，由 GPlates 网络服务（EarthByte）计算；高地数据取自 ETOPO1；地震数据取自美国地质调查局目录；板块边界取自 Bird（2003）；陆地数据取自 Natural Earth。
  \item 第~21 章（变质岩）：显微镜下的板劈理，照片由 Geoloshe 拍摄，CC0；锡罗斯岛的石榴云母片岩，照片由格雷厄姆·丘卡德拍摄，CC~BY~2.0；蓝晶石片岩和石英--黑云母片麻岩，照片由詹姆斯·圣约翰拍摄，CC~BY~2.0（均为维基共享资源）。
  \item 第~21 章：蓝晶石、红柱石和夕线石的稳定域取自 Holdaway 与 Mukhopadhyay（1993）。
  \item 第~22 章（俯冲带的岩浆作用）：1980 年 5 月 18 日圣海伦斯火山的喷发柱，照片由唐纳德·A.·斯旺森拍摄，USGS 提供，公有领域；花岗闪长岩、安山岩和内华达山脉，照片由詹姆斯·圣约翰拍摄，CC~BY~2.0（均为维基共享资源）。
  \item 第~22 章：地震数据取自 USGS 目录；火山取自史密森学会全球火山活动项目；板块边界取自 Bird（2003）；陆地数据取自 Natural Earth；干燥橄榄岩固相线取自 Hirschmann（2000），并采用 PREM 密度。
  \item 第~23 章（山脉的消亡）：国际空间站拍摄的阿巴拉契亚山脉，宇航员照片 ISS033-E-22378，NASA 提供，公有领域；石英砾岩，照片由詹姆斯·圣约翰拍摄，CC~BY~2.0（均为维基共享资源）。
  \item 第~23 章：剥蚀速率取自 Portenga 与 Bierman（2011）；密西西比河的泥沙负荷取自 Meade 与 Moody（2010）；衰变常数取自 Ahnert（1970）；瓦登丘陵是本系列虚构的地区。
  \item 第~24 章（地球的内部热量）：冰岛的内斯耶维利尔地热发电站，照片由格雷塔尔·伊瓦尔松拍摄，由 Fir0002 编辑，公有领域（维基共享资源）。
  \item 第~24 章：Davies 与 Davies（2010）的热流数据；放射性生热取自 KamLAND 地中微子测量（2011）以及四种同位素目前的放热量；岩石生热取自 Vilà 与 Fernández（2010）。
  \item 第~25 章（放射性测年与地球的年龄）：威廉·汤姆森（开尔文勋爵），1906 年，照片由伦敦 Russell and Sons 拍摄，史密森学会图书馆提供，无已知版权限制；卡尼恩迪亚布洛铁陨石，照片由杰弗里·诺特金拍摄，CC~BY-SA~2.5（均为维基共享资源）。
  \item 第~25 章：陨石铅同位素取自 Patterson（1956）；IUGS 和 IUPAC（2022）推荐的铷-87 衰变常数；历年估计值的历史取自 Holmes（1913、1927）、Patterson（1956）和美国地质调查局；瓦登丘陵是本系列虚构的地区。
  \item 第~26 章（早期地球及其大气）：达勒斯峡谷的条带状铁建造，照片由格雷厄姆·丘卡德拍摄，CC~BY~2.0；鲨鱼湾的叠层石，照片由保罗·哈里森拍摄，CC~BY-SA~3.0；塞多纳附近的红层，照片由宾夕法尼亚州的温迪拍摄，CC~BY-SA~2.0（均为维基共享资源）。
  \item 第~26 章：年轻太阳的光度取自 Gough（1981）；宙的界限与颜色取自国际年代地层图。
  \item 第~27 章（解读古气候）：南极冰芯中的气泡，照片由澳大利亚联邦科学与工业研究组织大气研究部门拍摄，CC~BY~3.0；浮游有孔虫，照片由玛西·罗宾逊拍摄，USGS 提供，公有领域；花粉粒，达特茅斯电子显微镜设施提供，公有领域（均为维基共享资源）。
  \item 第~27 章：EPICA 穹顶~C 温度（Jouzel 等人，2007）和南极冰芯二氧化碳（Bereiter 等人，2015），经 NOAA 古气候数据库提供；LR04 叠加曲线（Lisiecki 与 Raymo，2005）；ETOPO1 地形（NOAA）；NASA GISS（GISTEMP v4）；参照期取自 IPCC 第六次评估报告第一工作组第~2 章。
  \item 第~28 章（气候变化）：莫纳罗亚天文台；照片由布赖恩·瓦塞尔拍摄，NOAA 提供，公有领域（维基共享资源）。
  \item 第~28 章：查尔斯·戴维·基林与约翰·钦于 1997 年在莫纳罗亚；照片由约翰·米勒拍摄，NOAA 提供，公有领域（维基共享资源）。
  \item 第~28 章：数据来自 NOAA 全球监测实验室和斯克里普斯海洋研究所（二氧化碳）、NASA GISS（GISTEMP v4）、NOAA 卫星测高实验室（海平面）、Bereiter 等人（2015）经 NOAA 古气候数据库提供的冰芯数据，以及 IPCC 第六次评估报告第一工作组（强迫、归因与预测）。
\end{itemize}
